\documentclass[10pt]{article}
\usepackage[margin=0.72in]{geometry}
\usepackage{booktabs}
\usepackage{longtable}
\usepackage{array}
\usepackage{xcolor}
\usepackage[hidelinks]{hyperref}
\usepackage{microtype}

\newcommand{\artifact}[1]{\texttt{#1}}
\newcommand{\tbd}{\textbf{unverified hypothesis}}

\title{PegaFlow: Predicting When Shared KV Transport Beats Recompute on Heterogeneous Accelerators}
\author{Advisor-authored storyline draft; implementation and experiments remain student-owned}
\date{M0 paper-asset closure, August 2026}

\begin{document}
\maketitle

\begin{abstract}
External KV-cache storage can replace repeated prefill with storage lookup and
device transfer, but a cache hit is not itself a latency win.  PegaFlow already
provides a sidecar storage engine, a vLLM connector, host and SSD backing, and
cross-node transport.  Its repository also preserves both favorable and
unfavorable serving traces.  Those traces motivate a narrower systems question:
can a phase-aware break-even contract predict when saved prefill exceeds
packing, transfer, queueing, and materialization cost without reversing
time-to-first-token?  We formulate an architecture-guided mechanism hypothesis,
a fail-closed evidence ledger, and a matched evaluation contract.  Existing
artifacts remain exploratory because their historical code and runtime custody
is incomplete; consequently this draft reports no publishable performance
result.  On Ascend, the current contribution is a porting/adaptation carrier
plus a falsifiable native-mechanism hypothesis, not a completed
architecture-native mechanism.
\end{abstract}

\section{Introduction}

Long-context serving repeatedly materializes large KV states.  PegaFlow moves
that state outside an inference process and lets instances reuse it through a
shared store.  This separation creates a clean systems opportunity---replace
prefill with lookup and transfer---but it also adds a new critical path.  A hit
must be packed or registered, moved through host and device memory, scheduled
against concurrent transfers, materialized into the consumer layout, and
matched to the correct request lifecycle.  Average hit rate or aggregate
throughput therefore cannot determine whether reuse improves a request's SLO.

The repository's own evidence sharpens the problem.  Its preregistration states
the basic inequality ``saved prefill exceeds DMA cost.''  Historical traces
also retain regimes where burst concurrency reverses TTFT and where compressed
KV leaves little useful gap.  At the same time, repository issue~18 records a
custody defect: most historical PegaFlow and external runtime commits are not
reachable from durable refs.  We use these observations only to motivate a
question and freeze an evaluation; we do not promote their numbers into paper
claims.

This draft makes three bounded contributions.  First, it converts PegaFlow's
cache-hit story into a request-level break-even problem that accounts for the
entire transfer critical path.  Second, it states a phase-aware mechanism
hypothesis whose output is falsifiable by matched lifecycle receipts.  Third,
it defines the Ascend causal chain and the custody gate required before any
performance result is paper-admissible.  Mechanism implementation and all new
experiments remain the responsibility of the repository's student owners.

\section{The Seven-Question Research Contract}

\begin{longtable}{>{\raggedright\arraybackslash}p{0.05\linewidth}p{0.90\linewidth}}
\toprule
Q & Frozen answer \\
\midrule
1 & \textbf{Problem.} When a request can reuse shared KV state, which workload,
phase, and resource conditions predict that transfer will beat isolated
recompute without reversing TTFT?  The unit is a matched request lifecycle, not
an aggregate cache-hit ratio. \\
2 & \textbf{Importance.} A wrong decision adds packing, DMA, queueing, and
materialization to a request that recompute could have served sooner.  The
failure is user-visible tail latency and resource contention, even when average
throughput or utilization improves. \\
3 & \textbf{Hidden assumption and related-work boundary.} Storage caches,
generic KV connectors, and transfer engines commonly treat availability or a
local transfer estimate as sufficient.  PegaFlow asks whether that assumption
fails once request-level queue state, KV layout, accelerator copy semantics,
and prefill opportunity cost interact.  It is not a new cache policy, a general
RPC layer, or a claim that all external-KV systems lack observability. \\
4 & \textbf{Mechanism hypothesis.} A lifecycle-bound phase model estimates
$G=C_{\mathrm{prefill}}-(C_{\mathrm{pack}}+C_{\mathrm{queue}}+
C_{\mathrm{transport}}+C_{\mathrm{materialize}})$ and selects a fixed shared
path only when the lower confidence bound of $G$ clears a preregistered safety
margin; otherwise it retains the matched isolated path.  This selector and its
Ascend-native seam are \tbd. \\
5 & \textbf{Feasibility.} PegaFlow already exposes connector, prefetch, DMA,
queue, and request events and separates storage from the inference process.
That makes a minimal receipt-bound predictor implementable at an existing seam,
but feasibility is not evidence of accuracy or overhead. \\
6 & \textbf{Matched evaluation.} Compare the native isolated path, the existing
fixed PegaFlow path, and the proposed phase-aware treatment using identical
code/runtime pins, model and input order, service lifecycle, device allocation,
KV layout, concurrency, and graph mode.  Correctness must pass before TTFT,
TPOT, throughput, queue time, transfer time, and utilization are analyzed. \\
7 & \textbf{Takeaway.} A reusable result would be a predictive boundary for
when external KV transport is beneficial, together with a causal explanation
for latency--throughput reversal.  A negative result closes only the tested
phase model and seam. \\
\bottomrule
\end{longtable}

\section{Problem Model and Hidden Assumption}

For request $r$, let $P_r$ be the prefill work avoided by a reusable prefix and
let $X_r$ be the extra work introduced by external reuse.  A deployable choice
needs more than $P_r>X_r$ in expectation: it must preserve output identity,
respect the request and KV generation lifecycle, and meet a tail-SLO margin
under contention.  We decompose
\[
X_r = X_r^{pack}+X_r^{queue}+X_r^{transport}+X_r^{materialize}
      +X_r^{coordination}.
\]
The common hidden assumption is that a hit and a transfer-bandwidth estimate
adequately approximate $X_r$.  This drops shared queue delay, topology, block
layout conversion, graph-mode effects, and interference among copies.  It also
mixes producer warmup with consumer evidence unless the lifecycle is explicit.

PegaFlow's scope is deliberately narrower than routing or admission control.
The proposed decision changes only whether a reusable KV object follows the
already available shared path or the matched isolated path.  It does not choose
which request to serve, where to route it, how to compress its state, or which
tier should own it.

\section{Architecture Causal Chain for Ascend}

\paragraph{Architecture fact.}
The repository's Ascend carrier imports NPU memory across processes, uses
explicit host-pinned allocation and asynchronous device copies, and records a
NUMA-sensitive transfer path.  The serving carrier uses graph-mode execution
and a device-specific KV layout.  These are code and design facts; they do not
alone establish a performance contribution.

\paragraph{Invalidated assumption.}
A GPU-derived rule that treats a cache hit or nominal PCIe bandwidth as a
request-level win ignores explicit import/materialization, queue contention,
layout-dependent bytes, and graph-mode prefill cost.  The repository's preserved
negative examples make that assumption unsafe, but their exact numeric effects
remain exploratory until custody closes.

\paragraph{Real counterexample.}
Tracked serving-path traces include a burst regime in which the shared path is
worse and a compressed-KV regime near break-even.  They are real-path artifacts,
not simulation, yet issue~18 shows that most referenced code/runtime commits are
not durably retrievable.  We therefore use only the qualitative contradiction:
``hit implies win'' is false in at least one recorded real path.  No number from
those traces is a paper result in this draft.

\paragraph{Native mechanism.}
The candidate native mechanism combines lifecycle receipts with an
Ascend-aware phase estimate at the existing connector/transfer seam.  It must
account for copy-stream queue state, NUMA placement, materialized layout, and
graph-mode prefill opportunity cost.  No such reviewed mechanism and matched
result currently exist.  Accordingly, the present Ascend work is classified as
\textbf{porting/adaptation plus a pending native hypothesis}.

\paragraph{Predictive boundary.}
The shared path should win only when reusable prefill is large enough and copy
contention is low enough that the lower confidence bound of $G_r$ exceeds the
safety margin for the target SLO.  Short or heavily compressed prefixes,
bursty concurrent copies, topology mismatch, and receiver queueing predict the
opposite outcome.

\paragraph{Generalization class.}
The boundary may generalize to heterogeneous accelerators where external state
reuse substitutes device compute with host or fabric movement and explicit
materialization.  It does not automatically generalize to in-device prefix
caches or systems without lifecycle-visible transfer costs.

\section{Mechanism Hypothesis}

The proposed treatment has three conceptual pieces.  A \emph{receipt layer}
binds lookup, enqueue, transfer, materialization, and completion to one request,
KV identity, and generation.  A \emph{phase estimator} derives the expected
saved prefill and added critical-path components from matched, causally valid
observations.  A \emph{guarded decision} enables the shared path only when its
confidence-aware gain clears the frozen margin and correctness is already
qualified.

This is intentionally not an implementation specification.  Student owners
must choose the minimal carrier seam and demonstrate that instrumentation does
not dominate the estimated gain.  If the dominant delay is network or service
queueing outside PegaFlow's seam, issue~4 requires transferring the question to
the appropriate networking or scheduling project instead of expanding this
repository.

\section{Current Evidence Ledger}

Table~\ref{tab:evidence} separates execution class from paper admissibility.
``Real path'' never overrides missing provenance.

\begin{table}[ht]
\centering
\small
\begin{tabular}{>{\raggedright\arraybackslash}p{0.22\linewidth}
                >{\raggedright\arraybackslash}p{0.18\linewidth}
                >{\raggedright\arraybackslash}p{0.49\linewidth}}
\toprule
Artifact family & Evidence label & Paper boundary \\
\midrule
Trace-audit family & real-online trace, exploratory & Contains
request and transfer events, including negative examples.  Historical
code/runtime custody is incomplete; no numeric claim is admitted. \\
Performance T1--T5 family & real-online trace audit, exploratory & Some
audit manifests declare valid event conservation, while other runs fail closed.
Issue~18 blocks formal use until referenced PegaFlow, vLLM, and Ascend commits
and raw-log custody are durable. \\
Archived benchmark report & derived narrative & Historical
summary only; it is not raw evidence and cannot support a result. \\
Trace preregistration & contract & Defines matched arms,
oracles, and stop rules; it is not an experiment result. \\
Host and simulation tests & probe/simulation & Qualify local semantics only;
they cannot establish serving, NPU, or end-to-end effects. \\
Future matched matrix & \tbd & Requires durable carrier/runtime pins, raw
receipts, allocation identity, oracle closure, and an exact result manifest. \\
\bottomrule
\end{tabular}
\caption{Fail-closed evidence classification.}
\label{tab:evidence}
\end{table}

\section{Evaluation Contract}

\subsection{Arms and matching}

The strongest deployable baseline is the native isolated namespace under the
same runtime and graph mode.  The second baseline is fixed PegaFlow sharing with
no state-aware selection.  The treatment is the proposed phase-aware guarded
path.  Each comparison freezes the PegaFlow and runtime heads, official carrier
identity, dependency lock, model and configuration hash, prompt order, request
arrival trace, producer/consumer lifecycle, block layout, service flags, device
allocation, clock source, and fault/interference schedule.  Arm order alternates
and the unit of replication is an independent service lifecycle.

\subsection{Correctness oracle}

Every consumer request must match the isolated baseline's token IDs under the
repository's accepted tolerance.  Receipt conservation must report exactly one
eligible connector event, prefetch decision, transfer completion when required,
and materialization completion per request/KV generation.  Missing, duplicate,
orphan, cross-arm, cross-request, or fallback-only transfer evidence invalidates
the lifecycle.  Producer warmup is excluded from the consumer comparison.

\subsection{Metrics and preregistered gates}

Report full distributions of TTFT, TPOT, total latency, request throughput,
queue time, transfer time and bytes, cache-hit coverage, and accelerator/host
utilization.  The treatment may advance only if correctness is 100\%, every
claim-bearing lifecycle has complete receipts, and it beats both deployable
baselines without a preregistered unacceptable TTFT or TPOT regression.  Exact
numeric performance thresholds remain \tbd{} until owners freeze the target SLO
before observing the new matrix; this draft does not invent them.

\subsection{Stop conditions}

Stop after one bounded correctness-debug cycle if conservation or output
identity still fails.  Stop the selector hypothesis if its lower-bound decision
does not repeatedly separate winning and losing regimes, if oracle headroom
over the strongest fixed arm is absent, or if receipt overhead consumes the
predicted gain.  Stop the Ascend-native claim if only API adaptation, packaging,
or environment wiring differs from the generic path.  Negative outcomes close
only the tested model and seam; thresholds must not be moved after held-out
results are visible.

\section{Related-Work Boundary}

This repository does not yet contain a reviewed scholarly bibliography, so this
draft deliberately makes no external citation claim.  The comparison space is
nevertheless explicit.  External KV stores and connectors establish that KV
state can outlive one inference process.  Prefix caches establish reuse within
an engine.  Generic transfer engines optimize bytes, registration, and fabric
movement.  Queueing and SLO controllers choose request-level service actions.
PegaFlow's proposed paper boundary is the intersection: predicting a
request-level reuse break-even from lifecycle-bound transfer phases in an
existing sidecar data path.  A later related-work pass must add only verified
references and must narrow the novelty claim if any prior system already makes
the same receipt-bound prediction.

\section{Limitations and Two-Week Deliverable}

No result in this draft is paper-admissible.  Historical runtime refs are
incomplete, existing artifacts span multiple experimental revisions, and the
native Ascend mechanism is unimplemented.  The first week should close a
claim-bearing artifact inventory and durable code/runtime refs, then freeze the
minimal seam, receipts, baselines, and thresholds.  The second week should
produce host-level oracle tests and, only after independent allocation, one
matched real-path pilot.  The student owners perform all implementation,
environment setup, execution, debugging, and result generation.

\section{Conclusion}

PegaFlow already demonstrates a substantial external-KV system carrier, but its
research claim must be narrower than ``cache sharing is faster.''  The
falsifiable opportunity is to predict the request-level break-even between
saved prefill and the full transfer critical path.  Until provenance custody,
the native mechanism, and matched evaluation close, this remains a
well-specified hypothesis with exploratory evidence rather than a completed
systems result.

\section*{Bibliography status}
The accompanying \artifact{references.bib} is intentionally empty at this
stage.  No external citation was added without a repository-verifiable source.

\end{document}
